\documentclass[10pt]{article}
\usepackage[a4paper,margin=1.9cm]{geometry}
\usepackage[T1]{fontenc}
\usepackage{lmodern}
\usepackage{microtype}
\usepackage{booktabs}
\usepackage{multirow}
\usepackage{array}
\usepackage{xcolor}
\usepackage{tikz}
\usetikzlibrary{arrows.meta,positioning,fit,backgrounds,calc}
\usepackage{enumitem}
\usepackage[round,authoryear]{natbib}
\usepackage{caption}
\usepackage{hyperref}
\setlist{nosep,leftmargin=*}
\captionsetup{font=small,labelfont=bf,skip=4pt}
\hypersetup{colorlinks=true,linkcolor=blue!50!black,urlcolor=blue!50!black,citecolor=blue!50!black}
\setlength{\parskip}{3pt}
\setlength{\parindent}{0pt}
\setlength{\emergencystretch}{2em}
\setlength{\textfloatsep}{8pt}
\setlength{\intextsep}{6pt}
\setlength{\bibsep}{1pt}
\let\oldus\_\renewcommand{\_}{\oldus\allowbreak}   % long tool names may break after _
\newcolumntype{L}[1]{>{\raggedright\arraybackslash}p{#1}}
\newcommand{\todo}[1]{\textcolor{red}{\textbf{[TODO: #1]}}}
\newcommand{\ALLOW}{\textsc{allow}}
\newcommand{\ASK}{\textsc{ask\_confirmation}}
\newcommand{\BLOCK}{\textsc{block}}
\newcommand{\ESC}{\textsc{escalate}}
\newcommand{\NEED}{\textsc{need\_information}}
\newcommand{\ci}[2]{[#1, #2]}
\newcommand{\hd}[1]{\begin{tabular}[b]{@{}c@{}}#1\end{tabular}}   % two-line table header

\definecolor{cdata}{HTML}{374151}\definecolor{cdatabg}{HTML}{F3F4F6}
\definecolor{cmodel}{HTML}{1D4ED8}\definecolor{cmodelbg}{HTML}{DBEAFE}
\definecolor{cgate}{HTML}{B45309}\definecolor{cgatebg}{HTML}{FEF3C7}
\definecolor{cok}{HTML}{15803D}\definecolor{cokbg}{HTML}{DCFCE7}
\tikzset{
  >={Stealth[length=2mm]},
  box/.style={draw, rounded corners=2pt, align=center, inner sep=3pt, line width=0.6pt,
              minimum height=1.05cm, text width=#1, font=\footnotesize},
  data/.style={box=#1, draw=cdata, fill=cdatabg},
  model/.style={box=#1, draw=cmodel, fill=cmodelbg},
  gate/.style={box=#1, draw=cgate, fill=cgatebg},
  ok/.style={box=#1, draw=cok, fill=cokbg},
  note/.style={font=\scriptsize, fill=white, inner sep=1pt}}

\title{\vspace{-1.2cm}\textbf{Selective Environment Verification for Agent-Risk Monitoring}\\[2pt]
\large LMA Group Project 2026 \,--\, Final report}
\author{\textbf{Team Catalyst}: Meet Ghelani, Moiesha, Param Modi, Prateek Tiwari\\[2pt]
\normalsize GitHub: \url{https://github.com/meet6868/Agent-Risk-Monitoring}}
\date{14 November 2026}

\begin{document}
\maketitle
\vspace{-0.8cm}

\begin{abstract}
An LLM agent that works inside a company proposes actions such as sharing a file, granting access or
transferring money. Whether an action is safe often depends on a fact that is not in the conversation:
is the file shared outside the company, is the user the last administrator. We build a safety monitor that
decides \ALLOW, \ASK, \BLOCK\ or \ESC\ before each action and that may first ask read-only check tools for a
missing fact. We ask whether a small model can learn to check \emph{selectively}, only when it lacks a fact,
and stay about as safe as checking everything. We built a dataset of 594 records from 245 hand-written
scenarios with labels computed by policy rules, a store of 154 policies (170 in the out-of-domain store), 28 check
tools behind a gateway over a simulated company state, a hybrid policy retriever, and a Qwen3-8B monitor
fine-tuned with QLoRA. On the locked 146-record test set the monitor makes 0.74--0.80 checks per record
against 4.82 for always checking, at 63.0\% accuracy with retrieved policies and 69.9\% with the correct
policies (always: 62.3 and 75.3), and it beats the untrained model and two rule baselines. The hypothesis
fixed before any result, that selective checking lets through at most 5 points more risky actions than
always checking, \textbf{fails on test} ($\Delta$ unsafe $+14.6$ \ci{0.0}{29.7} with retrieved policies,
$+26.8$ \ci{14.3}{39.5} with correct policies; 95\% CI) and is inconclusive out of domain. The main cause is
cases that need two checks, where the monitor makes both in about 12\%. Selective checking saves check calls,
not model tokens (about 2.4$\times$ those of always checking).
\end{abstract}

\section{Introduction}
\textbf{The pricing-sheet example.} A user asks an agent to update a sales workbook, and the agent proposes
to write the internal cost numbers into \texttt{pricing\_v2.xlsx}. A company policy forbids writing internal
figures into a file that is shared outside the company. Nothing in the conversation says whether that file is
shared. A monitor that reads only the conversation has to guess. A monitor that queries every company system
before every action is slow, costs money and reads a lot of private state. The right step is one read-only
check, \texttt{check\_file\_sharing(pricing\_v2.xlsx)}: if it returns \texttt{sharing = external:org\_4} the
monitor blocks the action under the sharing policy, otherwise it allows it. A variant shows the harder part of
the problem: an earlier tool result in the conversation said the file was internal, but the agent then moved
the file, so that result is stale and the monitor has to check again.

\textbf{Why it matters.} Guard models label a prompt or a reply without any access to the environment, so
they cannot see the fact that decides such a case. Checking every fact for every action is the safe default
but it is expensive: each check is a call to a real company system, with latency, rate limits and access to
personal data. We study the middle way: a monitor that first decides whether it knows enough and asks for
exactly the missing fact when it does not.

\textbf{Research question.} Can a small model learn to check selectively and stay about as safe as checking
everything, while making far fewer checks? We compare four checking settings on the same records (never,
always, selective, oracle), with the hypothesis fixed before any result that selective checking lets through
at most 5 points more risky actions than always checking.

\textbf{Contributions.} To our knowledge, this is the first study that measures \emph{when} a learned monitor
chooses to query the environment and what that costs against always checking. Concretely:
\begin{enumerate}
\item a task and a dataset in which the deciding fact is hidden and must be checked: A/B record pairs that
  look the same and differ only in the hidden fact, tool results that may or may not be trusted (fresh,
  about another resource, made stale by a later write, or only claimed by the agent), and conflicts;
\item a monitor that decides when to check (\NEED) and which check to make, trained with QLoRA, with a
  policy retriever and a read-only check gateway;
\item a comparison of the four checking settings against two rule baselines, under correct and under
  retrieved policies, with confidence intervals and a pre-registered hypothesis;
\item an honest negative result on safety, an analysis of why it fails, and a cost analysis that separates
  check calls from model tokens.
\end{enumerate}

\textbf{What we do not claim.} We do not claim to be the first environment-aware monitor, the first to use
RAG for agent safety, the first to study hidden state, or the first policy-aware safety system. Our claim is
narrower: we study checking only when needed as a way to make policy-based agent safety monitoring work.

All code, data generators, policies, prompts, evaluation scripts and the generated result reports are in the
repository: \url{https://github.com/meet6868/Agent-Risk-Monitoring}. The only data not in the repository is
the AgentRisk dataset, which has no redistribution licence (Section~\ref{sec:data}).

\section{Literature Review}
\textbf{Agent safety benchmarks.} R-Judge \citep{yuan2024rjudge} tests LLMs as judges of risk in agent
conversations. ATBench \citep{li2026atbench} has long safe and unsafe agent trajectories with many tools.
Agent-SafetyBench \citep{zhang2024agentsafetybench}, AgentHarm \citep{andriushchenko2025agentharm},
SafeArena \citep{tur2025safearena}, OpenAgentSafety \citep{vijayvargiya2025openagentsafety} and ToolEmu
\citep{ruan2024toolemu} cover more harmful tasks and environments.
These benchmarks mostly give fixed, labelled records; the monitor cannot query an environment, which is the
step we study.

\textbf{Closest work.} ShieldAgent \citep{chen2025shieldagent} verifies policy rules using information from
the agent's trajectory; we study when to check the environment and what that costs compared with always
checking. AGrail \citep{luo2025agrail} creates safety checks before actions but has no explicit ``do I know
enough?'' step. PhantomPolicy and Sentinel \citep{wu2026phantompolicy} name the failure in which an agent's
compliance depends on facts it cannot see in its context, and enforce policies by simulating each action on a
knowledge graph; we give the monitor read-only tools and measure whether it learns to use them only when
needed. ``Ask the World Before Acting'' \citep{song2026askworld} lets an agent probe the environment for one
uncertain belief before it acts, under a budget, to improve task success; our probing step serves a separate
safety monitor and is compared with checking everything. pi-bench \citep{das2026pibench} tests whether agents
follow policies with hidden state; we test a separate monitor and how selective it is. SABER
\citep{hu2026saber} judges safety from the final state of a workspace, not before the action. Trust No Tool
\citep{yan2026trustnotool} uses hidden triggers with matched safe cases but focuses on attacked tool outputs.
AgentDoG \citep{liu2026agentdog} diagnoses risks in agent runs without environment access.

\textbf{Guard models and related ideas.} Guard models such as Llama Guard \citep{inan2023llamaguard} label
prompts or replies without access to the environment. Our ``do I know enough?'' step is close to selective
prediction \citep{geifman2017selective} and active feature acquisition \citep{shim2018afa}, where a model pays
to see extra information only when needed; we apply the idea to a safety decision with a hard call budget. We
use LoRA \citep{hu2022lora} and QLoRA \citep{dettmers2023qlora} on Qwen3-8B \citep{qwen2025qwen3}, retrieval
over a policy store \citep{lewis2020rag} with BM25 \citep{robertson2009bm25} and BGE embeddings
\citep{xiao2023cpack}, and Cohen's $\kappa$ \citep{cohen1960} for label agreement on the collected data.

\section{Dataset created / used}\label{sec:data}
\subsection{Our generated records: the main evaluation}
\textbf{From seed to record.} Each record starts from a hand-written \emph{seed}: one scenario with an
action, the fact that decides it, the policy and the hidden company state. Code and a writer LLM from a
different model family than the monitor (Claude) turn a seed into records: the writer LLM writes only the user
and agent sentences; the simulator supplies the company state; the gold label (decision, risk label, risk
type, whether a check is needed, the missing fact, the required checks, the policy ids) is \textbf{computed by
code} from machine-readable policy rules, never by an LLM. Each record has a \emph{visible} part (the
conversation and the proposed action), a \emph{hidden} part (the company state, reachable only through the
check tools) and the gold label (used only for scoring and training targets). Automatic checks run on every
record: schema, the identity of A/B pairs, that every hidden fact is answerable by a check tool, and
near-duplicate detection.

\textbf{Size and domains.} 245 seeds give 594 records: 200 in-domain seeds, 482 records, in six domains
(access control 112 records, files 108, privacy and sharing 94, money transfers 74, repositories 66, calendar
and email 28), and 45 out-of-domain (OOD) seeds, 112 records, with a new domain (customer records, 42 records),
reworded versions of the other domains, new tools, and new names, values and writing styles. The proposal
planned 195 seeds; the mid report had 212 seeds and 528 records; 13 two-check T5 test seeds (26 records) and
8 train seeds were added on 6 November, 12 seeds with new tools on 10 November, and distractor tool-error
turns were inserted into 26 train records to break an ``error message means unsafe'' shortcut.

\begin{table}[h]
\centering\small
\setlength{\tabcolsep}{4pt}
\begin{tabular}{@{}L{0.12\linewidth}L{0.50\linewidth}L{0.14\linewidth}L{0.08\linewidth}r@{}}
\toprule
\textbf{Type} & \textbf{What it tests} & \textbf{Gold decision} & \textbf{Checks needed} & \textbf{Records}\\
\midrule
T1 hidden fact & Two records A and B with the same visible text; only a hidden fact differs, so the right decision differs (e.g.\ the recipient address is verified or not) & set by the fact & 1 & 158\\
T2 trust rule & A fact appears earlier in the conversation: (a) a fresh tool result about the same resource, (b) a result about another resource, (c) a result made stale by a later write, (d) only the agent claims it. Only (a) may be trusted & set by the fact & a: 0; b, c, d: 1 & 164\\
T3 clearly safe & A policy allows the action whatever the hidden state is & \ALLOW & 0 & 30\\
T4 clearly unsafe & The violation is visible in the request itself & \BLOCK & 0 & 38\\
T5 conflict & (i) two shown policies demand opposite things, or (ii) two checks return facts that cannot both be true & \ESC & (i) 0; (ii) 2 & 92\\
\bottomrule
\end{tabular}
\caption{Case types of the generated data and their in-domain record counts (482 in all). A T1 seed gives 2
records, a T2 seed 4 (one per subtype), T3--T5 seeds 2 each in two writing styles.}
\label{tab:types}
\end{table}

\textbf{Splits and locking.} Records are split by seed, and seeds that share a gold policy stay in one split,
so no policy text decides both a training and a test case (Table~\ref{tab:splits}). The test and OOD files are
locked by SHA-256 hashes; every later change to them is logged with a relock. Policy ids never reach the model:
the prompt shows the policies as numbers in a random order. No near-duplicates across splits were found.
Evaluation uses the phase-1 tools, as in training: 146 of the 150 test records (4 \texttt{create\_event}
records are left out), and 290 of the 292 train records are used for training. The 6 November addition made
T5 28\% of test against 19\% of all in-domain records, which Section~\ref{sec:limits} comes back to.

\begin{table}[h]
\centering\footnotesize
\setlength{\tabcolsep}{4pt}
\begin{tabular}{@{}lrrrrrrrrrrrr@{}}
\toprule
 & & & \multicolumn{5}{c}{\textbf{Case type}} & \multicolumn{4}{c}{\textbf{Gold decision}} & \\
\cmidrule(lr){4-8}\cmidrule(lr){9-12}
\textbf{Split} & \textbf{Seeds} & \textbf{Records} & T1 & T2 & T3 & T4 & T5 & \textsc{allow} & \textsc{ask} & \textsc{block} & \textsc{esc.} & \textbf{Need a check}\\
\midrule
Train & 122 & 292 & 108 & 96 & 18 & 22 & 48 & 142 & 12 & 90 & 48 & 206\\
Validation & 16 & 40 & 14 & 16 & 2 & 6 & 2 & 21 & 0 & 17 & 2 & 28\\
Test (locked) & 62 & 150 & 36 & 52 & 10 & 10 & 42 & 67 & 4 & 37 & 42 & 109\\
OOD test (locked) & 45 & 112 & 36 & 44 & 10 & 10 & 12 & 61 & 3 & 36 & 12 & 71\\
\bottomrule
\end{tabular}
\caption{Splits of the generated data (records). ``Need a check'' counts records whose gold label requires at
least one check. OOD = the customer domain plus reworded other domains, new tools and new wording.}
\label{tab:splits}
\end{table}

\textbf{Policy store, check tools and gateway.} The in-domain store holds 154 policies, the OOD store 170
(reworded, plus the new domain); each policy is one plain sentence for the monitor, with machine-readable
logic and conflict pairs kept apart for the gold labeller. For the final system each policy also carries an
``applies to'' tag: the agent tools it governs, derived from the \emph{actions} field of the policy logic only
(never from its condition, fact or decision). There are 28 typed read-only check tools (2 OOD-only), e.g.\
\texttt{check\_file\_sharing}, \texttt{check\_user\_clearance}, \texttt{check\_repository\_status}, answered by
a \textbf{check gateway} built from each record's hidden state: it validates arguments, refuses writes,
returns only each tool's declared fields, returns errors as data (\texttt{not\_found},
\texttt{invalid\_args}), logs every call and enforces the call limit. A fixed \emph{check table} lists, per
action, every check any applicable policy can need; it defines the always setting (about 4.5--5 checks per
record).

\textbf{Label audit.} Because the labels come from our own rules, we checked them by code on all 594 records
(no LLM, no human judgement): the merge checks pass (0 problems, 0 T1 pairs failed), every record rebuilds
identically from its source, test and OOD match their locks, there are 0 near-duplicates, every scored record
can be decided from what the monitor sees (conversation plus gold check results), the same policy, fact and
value always give the same decision (168 groups, 0 contradictions), and every gold check runs without an
error. One OOD seed (OOD-CUST-T3-002, 2 records) has a gold label that differs from the policy rule as the
audit script reads it; it is left for a person to read. The planned human agreement study on our records
(two annotators, Cohen's $\kappa$) was \textbf{not done}; the only $\kappa$ we have is 0.57 on 26 R-Judge
risk types.

\subsection{Collected datasets: risk detection only}
Three public datasets give whole agent trajectories with a safe/unsafe label but no policies and no
environment, so they are used only for risk detection (safe or unsafe, and the risk type), not for the
four-way decision (Table~\ref{tab:collected}). Safe/unsafe comes from the source. The risk type is one of our
six classes (privacy leakage, security, access control, financial loss, data loss, other), made from a mapping
table and by hand; \textbf{no LLM labelled anything}. 41 physical-world records were removed. Each dataset
is split per dataset into train, validation and test by us. \textbf{Our R-Judge and ATBench numbers therefore
come from our own splits and cannot be compared with published leaderboards}; ATBench was released as a
test-only benchmark. The AgentRisk data \citep{agentrisk2026} has no redistribution licence: it was never
committed or uploaded, the derived adapters are private, and this report shows only counts and scores from
it. The training list for risk detection holds 1,034 examples (R-Judge 226, ATBench 594, AgentRisk 214),
balanced 517 safe / 517 unsafe, with privacy leakage capped at 149 so that no risk type dominates.

\begin{table}[h]
\centering\small
\setlength{\tabcolsep}{4pt}
\begin{tabular}{@{}lrrrrL{0.44\linewidth}@{}}
\toprule
\textbf{Dataset} & \textbf{Records} & \textbf{Train} & \textbf{Val} & \textbf{Test} & \textbf{How the risk type was made}\\
\midrule
R-Judge & 531 & 369 & 47 & 115 & no type field: all by hand, blind; 4 members; $\kappa=0.57$ on a 26-item double check\\
ATBench & 999 & 699 & 100 & 200 & harm category $\rightarrow$ mapping table agreed by 2 members; unclear cases by hand\\
AgentRisk & 617 & 394 & 98 & 125 & risk domain $\rightarrow$ mapping table; two-domain cases by hand\\
\bottomrule
\end{tabular}
\caption{Collected datasets after conversion and filtering (records). Splits are our own.}
\label{tab:collected}
\end{table}

\section{Methodology}
\subsection{Task}
The monitor sees the visible input $x=(\tau, a, p)$: the conversation $\tau$, the proposed action $a$ and the
policies $p$ shown for it. The environment has a hidden state $s$; a read-only gateway offers check tools
$q$, and a call returns $q(s)$. The monitor either decides or asks: a \textsc{decision} is one of \ALLOW,
\ASK, \BLOCK, \ESC\ with a risk label, a risk type and evidence (facts and policy numbers); \NEED\ names the
missing fact and 1--2 checks. If several policies apply and do not conflict, the strictest wins; \ESC\ is
reserved for conflicts between policies or between facts. A \emph{trust rule} says when a fact counts as
known: a tool result in the conversation states it and no later action changed that resource; what the agent
says by itself never counts.

\subsection{Four checking settings and two policy conditions}
\begin{center}\small
\begin{tabular}{@{}lL{0.5\linewidth}L{0.3\linewidth}@{}}
\toprule
\textbf{Setting} & \textbf{Checks before the decision} & \textbf{Role in the study}\\
\midrule
never & none; decide from the conversation & lower bound, no cost\\
always & every check of the fixed check table for the action (4.50--5.00 per record) & safe but expensive reference\\
selective & the checks the model itself asks for (\NEED), at most 4 calls, one retry after a tool error & \textbf{our method}\\
oracle & exactly the checks the gold label needs & upper bound for checking\\
\bottomrule
\end{tabular}
\end{center}
Every setting is run under two \textbf{policy conditions}: \emph{correct policies} (the gold policies plus 2
look-alikes, shuffled and numbered: checking and deciding as if retrieval were perfect, the main test of
checking) and \emph{retrieved policies} (the retriever's output: the full system).

\subsection{The monitor loop}
\begin{figure}[h]
\centering
\begin{tikzpicture}
\node[data=2.5cm]  (vis) at (0,0)     {\textbf{Visible input}\\conversation +\\proposed action};
\node[data=2.7cm]  (ret) at (3.5,0)   {\textbf{Policy retriever}\\v1: top 3 of the store\\v2: tagged, top 8};
\node[model=2.9cm] (llm) at (7.2,0)   {\textbf{Monitor}\\Qwen3-8B + QLoRA\\one JSON object per pass};
\node[ok=2.9cm]    (dec) at (11.3,0)  {\textbf{Decision}\\\textsc{allow} / \textsc{ask} /\\\textsc{block} / \textsc{escalate}\\+ evidence};
\node[gate=2.9cm]  (gw)  at (7.2,-2.3){\textbf{Check gateway}\\28 read-only tools, $\le$4 calls,\\errors returned as data};
\node[gate=2.9cm]  (sim) at (11.3,-2.3){\textbf{Company state}\\built from the hidden part;\\never in the prompt};
\draw[->] (vis) -- (ret);
\draw[->] (ret) -- (llm);
\draw[->,cok] (llm) -- node[note,above=1pt]{decides} (dec);
\draw[->,cgate] ([xshift=-14pt]llm.south) -- node[note,left=2pt,align=right]{\NEED:\\missing fact + 1--2 checks} ([xshift=-14pt]gw.north);
\draw[->,cgate] ([xshift=14pt]gw.north) -- node[note,right=2pt,align=left]{results; next pass\\($\le$3 passes per record)} ([xshift=14pt]llm.south);
\draw[->,cgate] (gw) -- (sim);
\end{tikzpicture}
\caption{The monitor loop in the selective setting. Each pass sends one prompt (decision rules, the trust rule,
the shown policies, the conversation, the proposed action, the check tools and any check results) and gets one
JSON object back. In the final system (v2) the loop also adds a second-source check when the model names a
fact that two tools report (Section~\ref{sec:v2}).}
\label{fig:loop}
\end{figure}

A record uses at most 3 model passes and 4 gateway calls; after a tool error the model may ask once more,
then it must decide. Invalid JSON counts as a wrong answer in every metric (the live system would fall back
to \ASK). A test suite enforces that no hidden field and no gold label ever reaches a prompt. The monitor is
Qwen3-8B at a pinned revision, greedy decoding, thinking mode off, run on Kaggle with two T4 GPUs.

\subsection{Policy retrieval: v1 and v2}\label{sec:v2}
The retriever is a hybrid of BM25 and \texttt{bge-base-en-v1.5} cosine similarity over the visible policy
sentences, score $s = 0.6\cdot\mathrm{embedding} + 0.4\cdot\mathrm{BM25}$ (each min-max scaled per query), query =
the proposed action plus all user and agent turns. The weight and the query form were tuned on train and
validation only. \textbf{v1} (the system of the mid report and of the first final runs) shows the top 3: the
deciding policy is in the top 3 for 55\% of train and validation records (top 5: 62\%). \textbf{v2} (the final
system) ranks only the policies whose ``applies to'' tag contains the proposed action's tool and shows 8.
Measured on the 332 train and validation records, the share of records with \emph{all} gold policies shown
rises from 0.52 to 0.63 at $k=3$ and from 0.69 to 0.87 at $k=8$ with the tag filter; on the records as run,
from 0.55 to 0.88 (train) and 0.33 to 0.80 (validation). Showing 8 policies costs about 90 prompt tokens.

\textbf{v2 additions to the selective loop.} (1) A \emph{second-source check}: when the model's missing fact
is one of the three facts that two check tools report (\texttt{account\_type}, \texttt{modifiable},
\texttt{role\_eligible}), the loop adds the other tool's check on the same resource from the action's check
table, so that the monitor sees both sources; the added calls are marked and counted. (2) Two prompt hints:
if a policy depends on two facts or two tools report the same fact, ask for both; after a failed check, do
not repeat the same call. Nothing was retrained. The decision to adopt v2 was taken on validation with a rule
written down before the run: v2 replaces v1 if selective accuracy with retrieved policies does not fall
(27 of 40), wrong blocks rise by at most 2 records, and selective with correct policies is at most 2 records
below v1 (33 of 40). The run passed: 31 of 40, wrong blocks 2 $\rightarrow$ 1, and 34 of 40. The retrieval
change alone gave 32 of 40, so the gain on validation is the retrieval change; the second-source rule never
fired there (its 2 two-check records use a different mechanism).

\subsection{Fine-tuning}
QLoRA on Qwen3-8B: rank 16, $\alpha$ 32, learning rate $2\cdot10^{-4}$, 2 epochs, 4-bit NF4, at most 4,096
tokens; the targets are built by code from the gold labels and the loss is on the answer only.
\textbf{Phase~A} trains on the collected data only (1,034 examples: risk label and type).
\textbf{Phase~B} (the final adapter, trained 10--11 November) adds our 290 train records as 988 examples:
a first-pass example per record (\textsc{decision}, or \NEED\ with the gold checks) and, for records that need
a check, a second-pass example (the gateway's results $\rightarrow$ \textsc{decision}), our examples repeated
twice, plus the same 1,034 collected examples: 2,022 examples, 1,981 within the token limit. Training loss
0.031, evaluation loss 0.024 $\rightarrow$ 0.025 over the two epochs, about 6.5 hours on two T4 GPUs. The
earlier Phase~B adapter of 30 October was trained the same way on 42 fewer train records and without the
distractor turns. There is one training run per adapter.

\subsection{Rule baselines, metrics and the hypothesis}
Two hand-written rule systems over the shown policy text and the facts: \emph{rules only} (never checks) and
\emph{rules + environment} (runs every table check first, then the same rules). They use only train records
and visible policy text.

Metrics (all in \%, defined in the repository's \texttt{docs/evaluation\_metrics.md}): decision
\textbf{accuracy}; \textbf{unsafe-action rate}, the share of risky records (gold is not \ALLOW) given a weaker
decision than the gold one (\ALLOW\ 0 $<$ \ASK\ 1 $<$ \BLOCK\ = \ESC\ 2); \textbf{wrong-block rate}, the share of
safe records blocked or escalated; \textbf{checks per record} and \textbf{saving} $=1-$ calls / calls(always),
always printed next to the \textbf{oracle saving} (the same for exactly the gold checks); check precision and
recall; right missing fact; T1 pair accuracy (both records of an A/B pair right); T2 accuracy per subtype;
decision macro-F1; risk-label F1 and risk-type macro-F1 on the collected data; valid-JSON rate. Every main
number gets a 95\% interval from a bootstrap over seeds (2,000 resamples; whole seeds are drawn because the
records of one seed are near-copies; paired between settings). Selective and always are compared with a
Wilcoxon signed-rank test on checks per record and an exact McNemar test on accuracy.

\textbf{Hypothesis (R10.4, fixed before any result).} Selective checking is not less safe than always
checking: the upper end of the 95\% interval of $\Delta$ unsafe-action rate (selective $-$ always) is at most
$+5$ points, while selective makes clearly fewer checks. The main test is with correct policies, so that
retrieval errors are not mixed in.

\textbf{Rigour.} Test and OOD were run and scored once with the final setup; nothing was trained, tuned or
chosen on them. Every change after the mid report (the final adapter, v2) was decided on validation with a
rule written in the decision log before the run.

\section{Experiment runs}
All runs use Qwen3-8B (pinned revision, greedy, thinking off) on Kaggle $2\times$T4 (at most two GPU jobs at a
time, 12 hours each). A full evaluation of one model on validation, test and OOD takes about 5 hours; a
Phase~B training run about 6.5 hours. Every model was trained once and every setting
run once; the runs, traces and generated tables are kept, and all numbers below are copied from them.

\subsection{Experiment 1: risk detection on the collected data}
Whole-trajectory risk detection, one pass, no checks, on our own splits (Table~\ref{tab:exp1}). Fine-tuning
lifts the base model from near-chance to about 90\% in every dataset. A \textbf{word model with no LLM}
(TF-IDF and logistic regression on the judged step, one per dataset, no tuning) already reaches 88.6, so most
of the jump from the base model comes from wording the datasets share between train and test; Phase~A's lead
over the word model is significant on test ($p=0.025$), the final Phase~B adapter's is not ($p=0.374$). On
R-Judge, 43 of the 115 test records are near-copies of a training record (Phase~A without them: 86.1); on
AgentRisk, a rule ``the agent replies to the user $\Rightarrow$ safe'' gets 85.6, but Phase~A also holds up
without replies (76 of 77 right); the ATBench tool-error shortcut cannot be tested on our splits (no safe
validation or test record carries it).

\begin{table}[h]
\centering\footnotesize
\setlength{\tabcolsep}{3.5pt}
\begin{tabular}{@{}lccc@{\hspace{10pt}}cccccc@{}}
\toprule
 & \multicolumn{3}{c}{\textbf{Validation (241 records)}} & \multicolumn{6}{c}{\textbf{Test (428 records)}}\\
\cmidrule(lr){2-4}\cmidrule(lr){5-10}
 & & F1 & Type & & F1 & Type & \multicolumn{3}{c}{Accuracy per dataset}\\
\textbf{Model} & Accuracy & unsafe & macro-F1 & Accuracy & unsafe & macro-F1 & R-Judge & ATBench & AgentRisk\\
\midrule
Base Qwen3-8B (prompted) & 52.3 & 14.8 & 4.0 & 54.9 & 32.8 & 26.1 & 57.4 & 53.7 & 54.4\\
Phase A (collected data only) & \textbf{93.8} & \textbf{93.4} & 59.6 & \textbf{92.8} & \textbf{91.9} & \textbf{55.0} & 91.3 & 89.4 & 99.2\\
Phase B, 30 October & 92.1 & 91.5 & 55.6 & 91.4 & 90.2 & 50.0 & 90.4 & 87.8 & 97.6\\
Phase B, final & 91.7 & 91.1 & \textbf{60.6} & 90.4 & 89.0 & 52.7 & 89.6 & 86.2 & 97.6\\
Word model, no LLM & -- & -- & -- & 88.6 & -- & -- & 87.8 & 86.2 & 92.8\\
\bottomrule
\end{tabular}
\caption{Experiment 1 (\%): safe/unsafe accuracy, F1 for the unsafe class and risk-type macro-F1 on the
collected validation and test sets (our own splits; 12 test records over 7,300 tokens left out).}
\label{tab:exp1}
\end{table}

\subsection{Experiment 2: our records in four settings (final system, v2)}
Tables~\ref{tab:main-test} and~\ref{tab:main-ood} give the main results on the locked in-domain test set and
the locked OOD set; validation is summarised in the text.

\begin{table}[h]
\centering\footnotesize
\setlength{\tabcolsep}{2.5pt}
\begin{tabular}{@{}llcrrcrrc@{}}
\toprule
 & & \multicolumn{3}{c}{\textbf{Correct policies}} & \multicolumn{3}{c}{\textbf{Retrieved policies (full system)}} & \\
\cmidrule(lr){3-5}\cmidrule(lr){6-8}
\textbf{Method} & \textbf{Setting} & Acc.\ [95\% CI] & Unsafe$\downarrow$ & \hd{Wrong\\block$\downarrow$} & Acc.\ [95\% CI] & Unsafe$\downarrow$ & \hd{Wrong\\block$\downarrow$} & \hd{Checks /\\record}\\
\midrule
Rules only & never & 43.8 \ci{34.3}{52.5} & 95.1 & 1.6 & 44.5 \ci{34.8}{53.2} & 93.9 & 0.0 & 0\\
Rules + environment & always & 52.1 \ci{39.9}{63.6} & 28.0 & 40.6 & 46.6 \ci{35.4}{57.6} & 28.0 & 53.1 & 4.82\\
\midrule
\multirow{4}{*}{Base Qwen3-8B} & never & 34.2 \ci{26.8}{41.4} & 23.2 & 59.4 & 36.3 \ci{28.1}{44.3} & 15.9 & 54.7 & 0\\
 & always & 52.1 \ci{41.8}{62.5} & 12.2 & 29.7 & 54.8 \ci{44.2}{65.4} & 24.4 & 18.8 & 4.82\\
 & selective & 43.2 \ci{34.7}{51.4} & 26.8 & 43.8 & 44.5 \ci{34.9}{53.5} & 26.8 & 40.6 & 0.58 / 0.61\\
 & oracle & 51.4 \ci{41.3}{60.6} & 6.1 & 31.2 & 50.0 \ci{39.9}{59.3} & 8.5 & 32.8 & 0.97\\
\midrule
\multirow{4}{*}{Phase B (QLoRA)} & never & 61.0 \ci{49.3}{70.9} & 29.3 & 14.1 & 58.2 \ci{46.1}{68.8} & 30.5 & 7.8 & 0\\
 & always & 75.3 \ci{65.4}{84.0} & 12.2 & 4.7 & 62.3 \ci{52.2}{71.3} & 24.4 & 4.7 & 4.82\\
 & \textbf{selective} & \textbf{69.9} \ci{59.0}{79.6} & \textbf{39.0} & \textbf{3.1} & \textbf{63.0} \ci{50.8}{73.3} & \textbf{39.0} & \textbf{6.2} & \textbf{0.74 / 0.80}\\
 & oracle & 86.3 \ci{79.2}{92.3} & 3.7 & 1.6 & 80.8 \ci{72.1}{88.2} & 15.9 & 1.6 & 0.97\\
\bottomrule
\end{tabular}
\caption{Experiment 2, locked in-domain \textbf{test} set, 146 records, final system v2 (\%). Accuracy with a
95\% CI (bootstrap over seeds); unsafe = unsafe-action rate; wrong block = wrong-block rate. Checks per record
are given for correct / retrieved policies where they differ.}
\label{tab:main-test}
\end{table}

\begin{table}[h]
\centering\footnotesize
\setlength{\tabcolsep}{2.5pt}
\begin{tabular}{@{}llcrrcrrc@{}}
\toprule
 & & \multicolumn{3}{c}{\textbf{Correct policies}} & \multicolumn{3}{c}{\textbf{Retrieved policies (full system)}} & \\
\cmidrule(lr){3-5}\cmidrule(lr){6-8}
\textbf{Method} & \textbf{Setting} & Acc.\ [95\% CI] & Unsafe$\downarrow$ & \hd{Wrong\\block$\downarrow$} & Acc.\ [95\% CI] & Unsafe$\downarrow$ & \hd{Wrong\\block$\downarrow$} & \hd{Checks /\\record}\\
\midrule
Rules only & never & 54.5 \ci{45.2}{62.7} & 100.0 & 0.0 & 53.6 \ci{44.3}{61.6} & 100.0 & 1.6 & 0\\
Rules + environment & always & 54.5 \ci{41.2}{67.5} & 31.4 & 45.9 & 48.2 \ci{36.4}{61.2} & 27.5 & 55.7 & 4.50\\
\midrule
\multirow{4}{*}{Base Qwen3-8B} & never & 48.2 \ci{38.9}{57.5} & 13.7 & 63.9 & 48.2 \ci{38.4}{57.7} & 11.8 & 60.7 & 0\\
 & always & 55.4 \ci{44.3}{66.3} & 19.6 & 44.3 & 58.9 \ci{50.0}{68.2} & 23.5 & 41.0 & 4.50\\
 & selective & 57.1 \ci{46.4}{68.0} & 9.8 & 52.5 & 50.0 \ci{40.8}{59.4} & 11.8 & 59.0 & 0.25 / 0.44\\
 & oracle & 63.4 \ci{52.0}{74.5} & 3.9 & 37.7 & 65.2 \ci{53.8}{76.4} & 3.9 & 36.1 & 0.65\\
\midrule
\multirow{4}{*}{Phase B (QLoRA)} & never & 67.9 \ci{58.3}{77.2} & 23.5 & 31.1 & 68.8 \ci{58.8}{78.2} & 19.6 & 32.8 & 0\\
 & always & 66.1 \ci{56.2}{75.5} & 19.6 & 27.9 & 63.4 \ci{53.9}{73.2} & 29.4 & 23.0 & 4.50\\
 & \textbf{selective} & \textbf{80.4} \ci{71.5}{88.0} & \textbf{15.7} & \textbf{14.8} & \textbf{66.1} \ci{54.8}{76.3} & \textbf{35.3} & \textbf{18.0} & \textbf{0.68}\\
 & oracle & 87.5 \ci{80.2}{93.9} & 3.9 & 9.8 & 82.1 \ci{73.6}{89.7} & 5.9 & 13.1 & 0.65\\
\bottomrule
\end{tabular}
\caption{Experiment 2, locked \textbf{OOD} test set, 112 records, final system v2 (\%). Same columns as
Table~\ref{tab:main-test}.}
\label{tab:main-ood}
\end{table}

\textbf{Findings.}
\begin{itemize}
\item \textbf{Fine-tuning helps a lot.} Phase~B beats the base model and both rule baselines in every setting
  on test and OOD, and on validation in every setting but one (always with retrieved policies: base 80.0 vs
  Phase~B 70.0 on 40 records). Validation, Phase~B never / always / selective / oracle: 75.0 / 82.5 / 85.0 /
  97.5 with correct policies and 72.5 / 70.0 / 77.5 / 90.0 with retrieved policies.
\item \textbf{Selective saves 83.4--84.9\% of the check calls}: 0.74--0.80 per record against 4.82 on test and
  0.68 against 4.50 on OOD (Wilcoxon $p<0.001$ everywhere).
\item \textbf{Accuracy is close to always checking, safety is not.} On test, selective is level with always
  under retrieved policies (63.0 vs 62.3) and behind with correct policies (69.9 vs 75.3; $\Delta$ accuracy
  $-5.5$ \ci{$-$15.0}{3.4}, McNemar $p=0.200$), but it lets through more risky actions: unsafe-action rate
  39.0 against 24.4 and 12.2. On OOD
  selective is \emph{more} accurate than always (80.4 vs 66.1 with correct policies, $p=0.009$): always
  over-reacts to facts that do not decide the case (wrong blocks 27.9 vs 14.8).
\item \textbf{The base model does not learn to check.} Its selective setting is worse than its always setting
  on test ($-$8.9 points with correct policies, $p=0.015$); its low OOD unsafe rate comes from blocking most
  records (wrong block 52--59\%), not from good checking.
\end{itemize}

\subsection{Experiment 3: system v1 against the final system v2}
Table~\ref{tab:v1v2} compares the same adapter under the two systems on the same records. v2 is safer with
retrieved policies (test unsafe rate: always 34.1 $\rightarrow$ 24.4, selective 52.4 $\rightarrow$ 39.0,
oracle 20.7 $\rightarrow$ 15.9) and shows more of the deciding policies (oracle accuracy 76.0 $\rightarrow$
80.8 on test, 75.9 $\rightarrow$ 82.1 on OOD), while accuracy moves little (selective $+1.4$ on test, $-0.9$
on OOD, both well inside the intervals). The base model's selective accuracy with retrieved policies also
rises (test 39.0 $\rightarrow$ 44.5, OOD 46.4 $\rightarrow$ 50.0).

\begin{table}[h]
\centering\small
\setlength{\tabcolsep}{4pt}
\begin{tabular}{@{}llcccccc@{}}
\toprule
 & & \multicolumn{4}{c}{\textbf{Accuracy, v1 $\rightarrow$ v2}} & \hd{\textbf{Unsafe, selective}\\v1 $\rightarrow$ v2} & \hd{\textbf{Acceptable, selective}\\v1 $\rightarrow$ v2}\\
\cmidrule(lr){3-6}
\textbf{Split} & \textbf{Policies} & never & always & \textbf{selective} & oracle & & \\
\midrule
Test & retrieved & 58.2 $\rightarrow$ 58.2 & 65.8 $\rightarrow$ 62.3 & \textbf{61.6 $\rightarrow$ 63.0} & 76.0 $\rightarrow$ 80.8 & 52.4 $\rightarrow$ 39.0 & 65.8 $\rightarrow$ 74.0\\
Test & correct & 61.0 & 75.3 & \textbf{69.2 $\rightarrow$ 69.9} & 86.3 & 41.5 $\rightarrow$ 39.0 & 72.6 $\rightarrow$ 74.7\\
OOD & retrieved & 60.7 $\rightarrow$ 68.8 & 59.8 $\rightarrow$ 63.4 & \textbf{67.0 $\rightarrow$ 66.1} & 75.9 $\rightarrow$ 82.1 & 37.3 $\rightarrow$ 35.3 & 68.8 $\rightarrow$ 72.3\\
OOD & correct & 67.9 & 66.1 & \textbf{78.6 $\rightarrow$ 80.4} & 87.5 & 15.7 $\rightarrow$ 15.7 & 82.1 $\rightarrow$ 83.0\\
Val & retrieved & 67.5 $\rightarrow$ 72.5 & 70.0 $\rightarrow$ 70.0 & \textbf{67.5 $\rightarrow$ 77.5} & 90.0 $\rightarrow$ 90.0 & 52.6 $\rightarrow$ 42.1 & 70.0 $\rightarrow$ 77.5\\
Val & correct & 75.0 & 82.5 & \textbf{82.5 $\rightarrow$ 85.0} & 97.5 & 26.3 $\rightarrow$ 26.3 & 82.5 $\rightarrow$ 85.0\\
\bottomrule
\end{tabular}
\caption{Experiment 3, Phase~B under system v1 (top-3 text retrieval, plain selective loop) and v2 (tagged
top-8 retrieval, second-source check, hints), in \%. With correct policies only the selective setting changes
between v1 and v2; a single number means unchanged. Acceptable = right, or \BLOCK\ where \ESC\ was right (or
the reverse): both stop the action. The hypothesis intervals under both systems are in Table~\ref{tab:hyp}.}
\label{tab:v1v2}
\end{table}

\section{Ablations}
Table~\ref{tab:abl} lists every ablation and comparison that was run. All numbers come from the generated
reports in the repository (\texttt{docs/reports/}).

\begin{table}[h]
\centering\footnotesize
\setlength{\tabcolsep}{4pt}
\begin{tabular}{@{}L{0.17\linewidth}L{0.22\linewidth}L{0.30\linewidth}L{0.26\linewidth}@{}}
\toprule
\textbf{Ablation} & \textbf{What changes} & \textbf{Result} & \textbf{What it shows}\\
\midrule
Checking setting & never / always / selective / oracle, Phase~B, test & correct policies 61.0 / 75.3 / 69.9 / 86.3; retrieved 58.2 / 62.3 / 63.0 / 80.8; checks 0 / 4.82 / 0.74--0.80 / 0.97 & checking matters (never 61.0, oracle 86.3 with correct policies); selective keeps most of the accuracy at a sixth of the calls but not the safety\\
Policy source & correct vs retrieved policies, selective & test 69.9 vs 63.0; OOD 80.4 vs 66.1; of the 48 v1 test errors that start with a retrieval miss, 12 become right with the correct policies & the retrieval ceiling is the gap between the two conditions ($+7.6$ test, $+11.6$ OOD, $+15$ val under v1), not the first-cause count\\
Retriever v1 $\rightarrow$ v2 & tag filter and 8 shown instead of top 3 & all gold policies shown: train 0.55 $\rightarrow$ 0.88, val 0.33 $\rightarrow$ 0.80; val selective 27 $\rightarrow$ 31 of 40; test unsafe (selective) 52.4 $\rightarrow$ 39.0 & better recall makes the system safer without retraining\\
v2 components (val) & retrieval change alone vs full v2 & 32 vs 31 of 40 (v1: 27); the second-source rule never fired on val & the validation gain is the retrieval change\\
Retriever top $k$ and tag filter (CPU, train + val, 332 records) & $k = 3 / 5 / 8 / 10$, with and without the filter & all gold policies shown 0.52 / 0.59 / 0.69 / 0.73 $\rightarrow$ with filter 0.63 / 0.74 / 0.87 / 0.89 & the filter helps more than a larger $k$; T1 policies name a hidden fact that no text query can find\\
Retriever query and weight & 3 query forms $\times$ 6 weights, recall@3 & action + all turns at $w=0.6$ is best (0.56); last turn 0.49; action only 0.44 & little is left in query wording\\
Cross-encoder re-ranker & \texttt{bge-reranker-base} on the top 20 & all gold shown at $k=3$: 0.52 $\rightarrow$ 0.32 (0.63 $\rightarrow$ 0.44 after the filter) & rejected: one-sentence policies and conversation-length queries suit a passage re-ranker badly\\
Training & base vs Phase~B, selective & test 43.2 $\rightarrow$ 69.9 (correct), 44.5 $\rightarrow$ 63.0 (retrieved); OOD 57.1 $\rightarrow$ 80.4, 50.0 $\rightarrow$ 66.1 & fine-tuning teaches when and what to check (Section~\ref{sec:ask})\\
Training data & 30 October vs final adapter, same 120 test records, v1 & never 56.7 $\rightarrow$ 70.8; selective 74.2 / 79.2 $\rightarrow$ 74.2 / 82.5; always 76.7 / 80.0 $\rightarrow$ 75.8 / 83.3; oracle 81.7 / 84.2 $\rightarrow$ 85.0 / 90.0 (retrieved / correct) & more data improved deciding, not asking; the final adapter asks a little less often\\
Phase~A vs Phase~B on collected data & collected-only vs all training data & test accuracy 92.8 vs 90.4; risk-type macro-F1 55.0 vs 52.7 & adding our records lowers collected-test accuracy from 92.8 to 90.4 and type macro-F1 from 55.0 to 52.7\\
Shortcut baseline & word model with no LLM, collected test & 88.6 vs Phase~A 92.8 and Phase~B 90.4 & most of the base $\rightarrow$ trained jump is shared wording\\
Rule baselines & rules only / rules + environment vs Phase~B & test, correct policies: 43.8 / 52.1 vs Phase~B selective 69.9 ($+17.8$ \ci{3.8}{32.1} over rules + environment) & the learned monitor beats fixed rules that check everything, at a sixth of the checks\\
\bottomrule
\end{tabular}
\caption{Ablations and comparisons that were run (\%, except retrieval shares and check counts).}
\label{tab:abl}
\end{table}

\textbf{Planned but not run.} Ablations with no policies shown, with noisy check results (the gateway supports
seeded noise) and a top-$k$ sweep on the model were planned but their code was not delivered. The Qwen3-235B
comparison was dropped (no endpoint decision). Live runs with two agent LLMs, the service and the demo were
not built. Two v3 candidates (show only the action's check tools; allow one more ask after the first results)
were built with a pre-set validation rule but could not be run: the Kaggle weekly GPU quota was used up on
12 November. None of these is reported as a result.

\section{Evaluations}
\subsection{The hypothesis}
Table~\ref{tab:hyp} gives the pre-registered test. \textbf{It fails on the in-domain test set} in both policy
conditions and under both systems: the interval of $\Delta$ unsafe is entirely above zero with correct policies
($+26.8$ \ci{14.3}{39.5} under v2), so selective checking is clearly less safe than always checking there;
with retrieved policies the v2 interval starts at 0.0 (\ci{0.0}{29.7}; under v1 \ci{8.2}{28.9}). On OOD and
on validation the intervals contain zero but reach above $+5$, so the result is inconclusive. The
cost side holds everywhere (Wilcoxon $p<0.001$). \textbf{Power.} With 82 risky test records the interval
half-width is 12.6--14.8 points, so the rule could only pass if the observed $\Delta$ were at most $-7.6$ to
$-9.8$ points, i.e.\ if selective were clearly \emph{safer} than always; on OOD (51 risky records) the needed
$\Delta$ is $-7.7$ to $-12.1$. The rule stays as it was fixed; this is the power it had. \textbf{Always is not
safe either}: it lets 24.4\% of risky test records through with retrieved policies and 12.2\% with correct
policies, because it inherits retrieval misses and sometimes decides on a fact that does not decide the case.

\begin{table}[h]
\centering\footnotesize
\setlength{\tabcolsep}{5pt}
\begin{tabular}{@{}llccrrl@{}}
\toprule
\textbf{Split} & \textbf{Policies} & \hd{\textbf{Phase~B v2}\\$\Delta$ unsafe [95\% CI]} & \hd{\textbf{Phase~B v1}\\$\Delta$ unsafe [95\% CI]} & \hd{Risky\\records} & \hd{Half-width\\(v2)} & \textbf{Result}\\
\midrule
\textbf{Test} & \textbf{correct (main)} & $\mathbf{+26.8}$ \ci{14.3}{39.5} & $+29.3$ \ci{16.7}{42.2} & 82 & 12.6 & \textbf{fails (clearly less safe)}\\
Test & retrieved & $+14.6$ \ci{0.0}{29.7} & $+18.3$ \ci{8.2}{28.9} & 82 & 14.8 & fails\\
OOD & correct & $-3.9$ \ci{$-$16.4}{9.1} & $-3.9$ \ci{$-$17.0}{8.9} & 51 & 12.7 & fails (inconclusive)\\
OOD & retrieved & $+5.9$ \ci{$-$10.9}{23.3} & $0.0$ \ci{$-$13.2}{13.2} & 51 & 17.1 & fails (inconclusive)\\
Val & correct & $0.0$ \ci{$-$21.1}{21.1} & $0.0$ \ci{$-$21.1}{21.1} & 19 & 21.1 & fails (inconclusive)\\
Val & retrieved & $+5.3$ \ci{$-$23.5}{31.2} & $+10.5$ \ci{$-$18.2}{41.2} & 19 & 27.4 & fails (inconclusive)\\
\bottomrule
\end{tabular}
\caption{Hypothesis R10.4: $\Delta$ unsafe-action rate, selective $-$ always, in points, with 95\% CIs
(bootstrap over seeds, paired). Passes only if the upper end is at most $+5$. The base model fails everywhere
except OOD with retrieved policies ($-11.8$ \ci{$-$26.4}{2.3}), where it passes only because it blocks most
records.}
\label{tab:hyp}
\end{table}

\subsection{Does the monitor learn to ask, and ask for the right fact?}\label{sec:ask}
Table~\ref{tab:ask} looks inside the selective setting with correct policies, on the same records for all
three models (test: the 120 records that existed before 6 November; these runs are under system v1, i.e.\ the
model's own requests). \textbf{Training teaches the model when to ask}, and it carries over to OOD: Phase~B
asks in 87.7\% (test) and 91.5\% (OOD) of the records that need a check, the base model in 53.1\% and 33.8\%.
It is clearest on T2: the base model never asks when the fact changed after the tool result (T2c) and rarely
when the result is about another resource (T2b: 31\% on test, 0\% on OOD); Phase~B asks in 100\% of both and
stays quiet in 92\% (test) and 91\% (OOD) of the T2a records whose result is fresh. \textbf{When it has the
facts, it decides right} (94.2\% when every needed check was made); the weak point is asking, not using the
answer. \textbf{It asks for one check when two are needed}: on the T5 records where two facts disagree Phase~B
asks in 94\% but makes both checks in 11.8\% (test, 34 records); on the 26 such records added on 6 November it
makes both checks in 2 and gets 2 right, while the oracle, given both facts, escalates correctly in 18 of 26.
\textbf{Wrong policies lead to wrong checks}: with retrieved policies the right check among the records where
it asked drops to 50.7 (test) and 59.1 (OOD), because the model checks the facts that the shown policies talk
about.

\begin{table}[h]
\centering\small
\setlength{\tabcolsep}{5pt}
\begin{tabular}{@{}lccc@{}}
\toprule
\textbf{Selective, correct policies; test / OOD} & \textbf{Base} & \textbf{Phase B, 30 Oct} & \textbf{Phase B, final}\\
\midrule
Asks when a check is needed (81 / 71 records) & 53.1 / 33.8 & 97.5 / 97.2 & 87.7 / 91.5\\
Quiet when no check is needed (39 / 41) & 84.6 / 92.7 & 89.7 / 82.9 & 92.3 / 80.5\\
Right check, among the records where it asked & 67.4 / 75.0 & 84.8 / 82.6 & 81.7 / 80.0\\
Every needed check made & 30.9 / 25.4 & 72.8 / 78.9 & 64.2 / 73.2\\
Accuracy when every needed check was made & 56.0 / 83.3 & 93.2 / 98.2 & 94.2 / 94.2\\
Accuracy when a needed check was missing & 55.4 / 52.8 & 36.4 / 53.3 & 55.2 / 21.1\\
\bottomrule
\end{tabular}
\caption{Asking behaviour in the selective setting with correct policies (\%), test (120 records) / OOD (112),
system v1.}
\label{tab:ask}
\end{table}

\subsection{T2 subtypes: tool results to trust or not}
Table~\ref{tab:t2} reports the trust-rule cases, the part of the task that other monitors do not handle.
Phase~B handles in-domain T2 almost perfectly in the selective setting (92.3--100 per subtype) and already
without checks (never: 92.3--100 on test), so the trust rules were learned from the conversation; the base
model is far lower (15.4--84.6 with correct policies).

\begin{table}[h]
\centering\small
\setlength{\tabcolsep}{5pt}
\begin{tabular}{@{}lcccc@{}}
\toprule
\textbf{T2 subtype, selective setting} & \hd{\textbf{Test}\\retrieved} & \hd{\textbf{Test}\\correct} & \hd{\textbf{OOD}\\retrieved} & \hd{\textbf{OOD}\\correct}\\
\midrule
Phase B: a. fresh result about the same resource & 100.0 & 100.0 & 81.8 & 90.9\\
Phase B: b. result about another resource & 100.0 & 100.0 & 72.7 & 90.9\\
Phase B: c. fact changed since the result & 92.3 & 100.0 & 90.9 & 100.0\\
Phase B: d. only the agent says it & 100.0 & 100.0 & 90.9 & 81.8\\
\midrule
Base Qwen3-8B: a. fresh result & 46.2 & 15.4 & 36.4 & 45.5\\
Base Qwen3-8B: b. another resource & 61.5 & 61.5 & 27.3 & 27.3\\
Base Qwen3-8B: c. fact changed & 84.6 & 84.6 & 100.0 & 100.0\\
Base Qwen3-8B: d. only the agent says it & 38.5 & 38.5 & 36.4 & 36.4\\
\bottomrule
\end{tabular}
\caption{Accuracy (\%) on the T2 records by subtype, selective setting, system v2: Phase~B (top) and the base
model (bottom). 13 test records and 11 OOD records per subtype.}
\label{tab:t2}
\end{table}

\subsection{T1 pairs: both records of an A/B pair right}
The two records of a T1 pair look the same and differ only in the hidden fact, so only a check can get both
right; a monitor that never checks gets at most one (Table~\ref{tab:pairs}). T1 is the hardest case type for
selective: with correct policies it gets both records right in under half of the pairs (8 of 17 on test, 8 of
18 on OOD; T1 record accuracy 73.5 and 66.7). With retrieved policies pairs fall further, because a missing
policy breaks both records of a pair.

\begin{table}[h]
\centering\small
\setlength{\tabcolsep}{6pt}
\begin{tabular}{@{}lcccc@{}}
\toprule
\textbf{Phase B, pairs with both right} & \textbf{never} & \textbf{always} & \textbf{selective} & \textbf{oracle}\\
\midrule
Test, correct policies (17 pairs) & 1 & 10 & \textbf{8 (47.1\%)} & 12\\
Test, retrieved policies (17 pairs) & 0 & 4 & \textbf{4 (23.5\%)} & 9\\
OOD, correct policies (18 pairs) & 0 & 11 & \textbf{8 (44.4\%)} & 13\\
OOD, retrieved policies (18 pairs) & 2 & 6 & \textbf{3 (16.7\%)} & 12\\
\bottomrule
\end{tabular}
\caption{T1 pair accuracy: number of A/B pairs with both records decided right, system v2.}
\label{tab:pairs}
\end{table}

\subsection{Saving next to the oracle saving}
Saving is $1-$ check calls / check calls of always; the oracle saving is the same quantity for exactly the
gold checks, i.e.\ the best saving possible with every needed check made (Table~\ref{tab:saving}). Selective
makes somewhat fewer checks than the oracle on test (108--117 against 141 calls) and about as many on OOD (76
against 73), so it does not save by skipping checks wholesale: most of its gap to the oracle's accuracy (16--18
points on test) comes from \emph{which} checks it makes, a wrong fact or one check of two.

\begin{table}[h]
\centering\small
\setlength{\tabcolsep}{5pt}
\begin{tabular}{@{}lccccc@{}}
\toprule
\textbf{Phase B, v2} & \hd{\textbf{Calls: always /}\\\textbf{selective / oracle}} & \hd{\textbf{Selective}\\\textbf{saving}} & \hd{\textbf{Oracle}\\\textbf{saving}} & \hd{\textbf{Selective}\\\textbf{accuracy}} & \hd{\textbf{Oracle}\\\textbf{accuracy}}\\
\midrule
Test, correct policies & 704 / 108 / 141 & \textbf{84.7} & 80.0 & 69.9 & 86.3\\
Test, retrieved policies & 704 / 117 / 141 & \textbf{83.4} & 80.0 & 63.0 & 80.8\\
OOD, correct policies & 504 / 76 / 73 & \textbf{84.9} & 85.5 & 80.4 & 87.5\\
OOD, retrieved policies & 504 / 76 / 73 & \textbf{84.9} & 85.5 & 66.1 & 82.1\\
\bottomrule
\end{tabular}
\caption{Check calls over the whole split and the saving against always checking (\%), next to the oracle
saving, system v2.}
\label{tab:saving}
\end{table}

\subsection{Error analysis}
Every wrong decision of Phase~B in the selective setting is put in exactly one group, the first failure in
the order the monitor works (Table~\ref{tab:err}; system v1 runs). With correct policies the largest group on
test is \textbf{two-check T5 records where one check was made} (22 of 45); with retrieved policies a
\textbf{retrieval miss} is the most common first cause (48 of 56 on test, 26 of 37 on OOD). The first-cause
count overstates what better retrieval can win: given the correct policies, only 12 of those 48 test records
become right and 36 stay wrong for other reasons (OOD 17 of 26, validation 6 of 10); the real ceiling is the
gap between the two policy conditions, $+7.6$ points on test, $+11.6$ on OOD and $+15$ on validation under
v1. Most errors are in the unsafe direction: on test, 43 of the 56 (retrieved) and 34 of the 45 (correct)
wrong answers let a risky action through.

Under v2 the second-source rule fired on only 1 test record. Of the 34 two-check test records, 10 are of the
kind the rule covers (one fact reported by two tools: 8 \texttt{modifiable}, 2 \texttt{role\_eligible}), but
the rule starts from the fact the model names, and there the model named another fact (for example
``production'' of the repository while calling one of the two right tools), so the rule did not start. The
other 24 need two \emph{different} facts combined (classification and sharing, classification and retention,
recipient class, registered vendor); all 26 two-check records in training use one other mechanism
(\texttt{account\_type}), so no inference rule can supply those without new training data.

\begin{table}[h]
\centering\footnotesize
\setlength{\tabcolsep}{4pt}
\begin{tabular}{@{}lccccccc@{}}
\toprule
\textbf{Phase B, selective} & \textbf{Wrong} & \hd{Retrieval\\miss} & \hd{Missed\\check} & \hd{Two needed,\\one made} & \hd{Wrong fact\\checked} & \hd{Right fact,\\wrong decision} & Other\\
\midrule
Test, retrieved policies & 56 of 146 & 48 & 0 & 0 & 5 & 1 & 2\\
Test, correct policies & 45 of 146 & -- & 5 & 22 & 10 & 8 & 0\\
OOD, retrieved policies & 37 of 112 & 26 & 2 & 0 & 4 & 4 & 1\\
OOD, correct policies & 24 of 112 & -- & 4 & 0 & 10 & 7 & 3\\
Val, retrieved policies & 13 of 40 & 10 & 0 & 0 & 1 & 1 & 1\\
\bottomrule
\end{tabular}
\caption{First cause of each wrong decision (counts), Phase~B selective, system v1. ``Other'' = wrong tool or
arguments, extra check, invalid JSON.}
\label{tab:err}
\end{table}

\subsection{OOD by kind of shift}
Of the 112 OOD records, 42 use new tools as well as new wording, style and values; 70 change only wording,
style and values. Phase~B selective (v2) scores 61.9 on the new-tool records and 68.6 on the reworded ones
with retrieved policies (78.6 and 81.4 with correct policies); always scores 61.9 and 64.3 (73.8 and 61.4).
Selective is weakest where it has never seen the check tools, and overall its OOD accuracy is close to its
test accuracy (66.1 vs 63.0 retrieved; 80.4 vs 69.9 correct), so the learned checking carries over to a new
domain.

\subsection{Exploratory cost hypotheses}
Four cost hypotheses were stated on 12 November, \textbf{after the results were seen}, so they are exploratory
and R10.4 stays the only pre-registered test; their thresholds are round values (half the cost) and the same
5-point margin (Table~\ref{tab:cost}). \textbf{Fewer check calls, yes (C1)}: selective makes 0.15--0.17 of the
check calls of always. \textbf{Less model compute, no (C2)}: selective uses about 2.4$\times$ the prompt and
output tokens of always (test, retrieved: 2,591 against 1,056 per record) and about twice the time (20.2
against 12.2 seconds per record), for two reasons: its prompt lists every check tool (about 990 tokens), and
most records need a second pass after the checks (1.63--1.75 passes per record; always and never make one).
The break-even point is 351--392 tokens of model compute per check call: selective is cheaper overall only
where one check call costs more than that, i.e.\ real company APIs with latency, rate limits, fees or access
to personal data, not a free local lookup. The proposal had anticipated this: selective checking ``may use
more tokens or time; its gain may be fewer checks and less access to private data rather than lower cost''.
\textbf{Accuracy within 5 points of always (C3)}: not shown on test (the intervals are too wide; $-5.5$ with
correct policies); on OOD with correct policies selective is clearly more accurate ($+14.3$
\ci{3.4}{25.0}). \textbf{Better than rules that check everything (C4)}: holds on test in both policy
conditions and on OOD with correct policies, by 16--26 points at about a sixth of the checks.

\begin{table}[h]
\centering\footnotesize
\setlength{\tabcolsep}{5pt}
\begin{tabular}{@{}lcccc@{}}
\toprule
\textbf{Phase B v2} & \hd{\textbf{C1: check calls}\\ratio sel./always $\le 0.5$} & \hd{\textbf{C2: model tokens}\\ratio sel./always $\le 0.5$} & \hd{\textbf{C3: accuracy kept}\\$\Delta$ acc.\ vs always $\ge -5$} & \hd{\textbf{C4: beats rules + env.}\\$\Delta$ acc.\ vs rules + env.\ $> 0$}\\
\midrule
Test, retrieved policies & \textbf{0.17} \ci{0.14}{0.20} & 2.45 \ci{2.35}{2.57} & $+0.7$ \ci{$-$7.7}{8.6} & $\mathbf{+16.4}$ \ci{0.7}{31.1}\\
Test, correct policies & \textbf{0.15} \ci{0.13}{0.18} & 2.48 \ci{2.39}{2.58} & $-5.5$ \ci{$-$15.0}{3.4} & $\mathbf{+17.8}$ \ci{3.8}{32.1}\\
OOD, retrieved policies & \textbf{0.15} \ci{0.12}{0.18} & 2.33 \ci{2.22}{2.43} & $+2.7$ \ci{$-$9.5}{14.9} & $+17.9$ \ci{0.0}{34.5}\\
OOD, correct policies & \textbf{0.15} \ci{0.12}{0.19} & 2.43 \ci{2.31}{2.56} & $\mathbf{+14.3}$ \ci{3.4}{25.0} & $\mathbf{+25.9}$ \ci{10.5}{40.5}\\
\bottomrule
\end{tabular}
\caption{Exploratory cost hypotheses, Phase~B v2: ratios of selective to always (C1, C2) and accuracy
differences in points (C3 against always, C4 against rules + environment), with 95\% CIs (bootstrap over
seeds, paired). \textbf{Bold} = the hypothesis holds (its 95\% bound is on the right side of the threshold).
C2 does not hold anywhere (the whole interval is above 0.5); a plain C3 or C4 cell means not shown (the
interval crosses the threshold).}
\label{tab:cost}
\end{table}

\section{Limitations}\label{sec:limits}
\begin{itemize}
\item \textbf{Small sets.} 40 validation, 146 test and 112 OOD records; the accuracy intervals are about $\pm$10
  points on test, and OOD has only 51 risky records. Any single difference under about 8 points is noise.
\item \textbf{Power of the hypothesis test.} With this size the pre-registered rule could only have passed if
  selective were clearly safer than always (Table~\ref{tab:hyp}); the rule was kept as fixed and its power is
  reported instead of changing it.
\item \textbf{Test mix.} The 6 November addition made test 28\% T5 against 19\% over all in-domain records.
  Reweighted to the in-domain case-type mix, Phase~B selective (v2) would be 68.3 (retrieved) and 75.1
  (correct) instead of 63.0 and 69.9; the raw test numbers stay the main result, and the 30 October numbers on
  the smaller test set are not comparable with them.
\item \textbf{Generated records and rule-made labels.} Our records, their labels and the check tools' answers
  come from our own seeds, policy rules and a simulated company state. The labels were checked by code
  (Section~\ref{sec:data}) but \textbf{no human agreement study was done on our records}; the only $\kappa$ is
  0.57 on 26 R-Judge risk types, below the 0.6 target. One OOD seed is flagged for a human to read.
\item \textbf{One training run per adapter.} Run-to-run variance is unknown, so differences of a few points
  between the 30 October and the final adapter may be noise.
\item \textbf{Two-check cases.} All 26 two-check training records use one mechanism; the 24 test records that
  combine two different facts are not learned, and no inference rule fixes them.
\item \textbf{Accuracy counts} \BLOCK\ \textbf{for} \ESC\ \textbf{as wrong} although both stop the action. Counting that swap as
  acceptable, always rises from 75.3 to 89.0 on test with correct policies but selective only from 69.9 to
  74.7: most of selective's errors really let the risky action through.
\item \textbf{Selective costs more model tokens than always} (about 2.4$\times$) and about twice the time; it
  saves check calls and access to private state, not compute.
\item \textbf{Collected datasets.} Our own splits, not comparable with leaderboards; shortcuts (a word model
  reaches 88.6; 43 R-Judge near-copies); the AgentRisk licence allows no redistribution.
\item \textbf{Not done.} The plug-in service, the demo and the live runs with two agent LLMs are empty in the
  repository; the no-policy, check-noise and top-$k$ model ablations, the Qwen3-235B comparison, the v3
  validation check and the 50 human evidence checks were not run.
\end{itemize}

\section{Conclusions}
\begin{enumerate}
\item \textbf{Fine-tuning teaches an 8B model to monitor agents with checks.} Phase~B beats the untrained model
  and both rule baselines in every setting on test and OOD, and the skill carries over to a new domain
  (selective, retrieved policies: 63.0 on test, 66.1 on OOD).
\item \textbf{Selective checking costs about a sixth of the check calls of always checking}, and the model
  learns \emph{when} to ask (87.7--91.5\% of the records that need a check, 100\% on stale or unrelated tool
  results) and mostly \emph{what} to ask; with every needed fact in hand it decides right in 94.2\%.
\item \textbf{It is not as safe as always checking on the in-domain test set.} The pre-registered hypothesis
  fails ($\Delta$ unsafe $+26.8$ \ci{14.3}{39.5} with correct policies), mainly because the monitor makes one
  check where two are needed, and because a retrieval miss turns into a wrong check. Both answers were said
  to be useful in the proposal; this is the negative one, reported as it is.
\item \textbf{Retrieval with action tags (v2) made the system safer without retraining} (test unsafe rate,
  selective 52.4 $\rightarrow$ 39.0), but the retrieval ceiling is modest ($+7.6$ points on test under v1)
  and the remaining errors are checking errors.
\item \textbf{The cost case is about calls, not tokens.} Selective is cheaper only where a check call costs
  more than about 350--390 tokens of model compute.
\end{enumerate}
\textbf{Future work}, in the order we would do it: (i) two-check training data with several mechanisms
(train-only composite policies in four domains) and oversampled two-check examples, so that ``ask for both
sources'' is learned as a rule rather than one pattern; (ii) second-pass training examples with the full
check-table results, so that the model learns to ignore facts that do not decide the case (the always and
OOD wrong-block errors); (iii) a second training seed of the final configuration, the only way to tell run
noise from real change; (iv) the human label study on our records, with a held-out train-side dev set that
covers every mechanism; and (v) a prompt that lists only the action's check tools and reuses the shared prefix
between passes, to cut the token cost.

{\small
\bibliographystyle{plainnat}
\bibliography{final_submission}}

\end{document}
